\subsection{The politico-economic equilibrium with CRRA}\label{\refsec:PEE}
With CRRA preferences, the past savings level $s_{t-1}$ becomes a relevant state variable for the political problem: the marginal indirect utilities are no longer pure log-derivatives but carry consumption \emph{levels}, and levels depend on the inherited stock. What is identified at each $t$ is therefore a policy \emph{function} $\tau_t(s_{t-1})$ rather than a number, built by backward recursion. The state-space reduction survives otherwise intact --- the savings \emph{shares} $s_{t-1,i}/s_{t-1}$ remain closed form in $\tau_t$, so the single aggregate level is the only state --- and the machinery of \cref{\refsec:RobustRoot} applies along the policy dimension unchanged. What is new is the grid structure over the state, and a fixed-point problem that pins the endogenous state down.

\smalltitle{Grids.} Let $\mathcal{S} = \left\lbrace s^{(1)}, ..., s^{(M_s)} \right\rbrace$ with $s^{(1)}$ = $l_s>0$ and $s^{(M_s)} = u_s>l_s$ denote a grid of states. Realistic ranges come from steady state dynamics: $s^*(\tau)$ is monotonically decreasing in $\tau$ with $s^* \rightarrow 0$ as $\tau \rightarrow 1$, so $l_s$ is a low value serving numerical stability and the upper bound can be taken as e.g.\ $u_s = 1.25 \times s^*(0)$. Let $\mathcal{S\tilde{T}}$ denote the Cartesian product of the state grid and the extended policy grid.

For $t<T$ we additionally need a grid of \emph{candidate} values for the endogenous state $s_t$ itself; call it $\mathcal{S}'$. It plays a different role from $\mathcal{S}$ --- $\mathcal{S}$ indexes the predetermined states we solve \emph{at}, while $\mathcal{S}'$ is searched \emph{over} to locate an equilibrium --- so although the two may hold the same values, it is worth keeping them separate, not least because the accuracy of the search below is governed by $\mathcal{S}'$ alone and may warrant a finer spacing.

\smalltitle{The economic equilibrium at $t<T$ is a fixed point in $s_t$.}
Suppose we have solved period $t+1$ and hold continuation policies $\tau^{t+1}(\cdot)$ and $h^{t+1}(\cdot)$, both functions of the state $s_t$ entering $t+1$. For a candidate current tax $\tau_t$ and a candidate state $s_t$, the auxiliary equations \eqref{\refeq:auxiliary} can be evaluated in one forward pass:
\begin{subequations}\label{\refeq:stateApprox}
\begin{align}
    \bm{B}_{t+1} &= \bm{B}_{t+1}\left(s_t,\, h^{t+1}(s_t)\right), \label{\refeq:stateApprox:B}\\
    \Gamma_{s,t} &= \Gamma_{s,t}\left(\bm{B}_{t+1},\, \tau^{t+1}(s_t)\right), \label{\refeq:stateApprox:Gammas}\\
    \Theta_{h,t} &= \Theta_{h,t}\left(\tau_t,\, \tau^{t+1}(s_t),\, \theta_{t+1},\, \Gamma_{s,t}\right), \label{\refeq:stateApprox:Thetah}\\
    \Theta_{s,t} &= \Theta_{s,t}\left(\Theta_{h,t},\, \Gamma_{s,t}\right). \label{\refeq:stateApprox:Thetas}
\end{align}
\end{subequations}
The candidate $s_t$ is an equilibrium at $s_{t-1}$ precisely when it reproduces itself through \eqref{\refeq:auxiliary:s}. Defining the residual
\begin{align}\label{\refeq:stateResidual}
    \mathcal{R}_t\left(\tau_t, s_t;\, s_{t-1}\right) \equiv \Theta_{s,t}\left(\tau_t, s_t\right)\left(\frac{s_{t-1}}{\nu_t}\right)^{\frac{\alpha(1+\xi)}{1+\alpha\xi}}-s_t,
\end{align}
the state policy function $s_t(\tau_t, s_{t-1})$ is the root of \eqref{\refeq:stateResidual} in its second argument. This is a genuine root-finding problem in $s_t$ --- not a recursion we can unroll --- because $s_t$ enters the right-hand side through $\bm{B}_{t+1}$ and the continuation policies. Note that it is a \emph{root} problem, not a maximisation: any crossing is a solution, so the direction-filtering and objective-integration apparatus of \eqref{\refeq:candidates} is not used here.

\smalltitle{The fixed point is two-dimensional, not three.}
Read \eqref{\refeq:stateApprox} again and note that \emph{none} of $\bm{B}_{t+1}$, $\Gamma_{s,t}$, $\Theta_{h,t}$, $\Theta_{s,t}$ depends on $s_{t-1}$: the predetermined state enters \eqref{\refeq:stateResidual} only through the explicit factor $(s_{t-1}/\nu_t)^{\alpha(1+\xi)/(1+\alpha\xi)}$. The costly part of the state approximation is therefore a function of $(\tau_t, s_t)$ alone, evaluated once on $\tilde{\mathcal{T}}\times\mathcal{S}'$, after which the residual for \emph{every} $s_{t-1}\in\mathcal{S}$ follows by multiplying $\Theta_{s,t}$ by a scalar and subtracting $s_t$. What looks like a three-dimensional problem thus costs a two-dimensional one plus a broadcast, which is what keeps the approach affordable as $\mathcal{S}$ is refined.

\smalltitle{Feasibility.}
The root of \eqref{\refeq:stateResidual} need not lie inside $\mathcal{S}'$. For a given $s_{t-1}$, large $\tau_t$ depresses savings enough that the implied $s_t$ falls below $l_s$, and conversely at small $\tau_t$; outside that window no equilibrium is located and every object downstream of $s_t$ is undefined. We therefore carry an explicit feasibility mask over $\tilde{\mathcal{T}}\times\mathcal{S}$, marking the pairs $(\tau_t, s_{t-1})$ for which a root was found, and require at least two feasible $\tau_t$ per state before attempting to select a maximum there. Since $s_t(\tau_t)$ is decreasing in $\tau_t$, the feasible set is expected to be a contiguous interval in $\tau_t$, but nothing relies on that: masking infeasible cells individually is no harder and does not silently mis-handle a non-monotone case.

\smalltitle{Which derivatives may be taken on the grid.}
Once $s_t(\tau_t,s_{t-1})$ is known, all four marginal indirect utilities take the form (level)$^{1-1/\rho}\times$(log-derivative):
\begin{subequations}\label{\refeq:PEEcrraTerms}
\begin{align}
    \frac{d\upsilon_{1,t}^i}{d\tau_t} &= \left(\hat{c}_{1,t}^i\right)^{1-1/\rho}\frac{d\ln\left(\hat{c}_{1,t}^i\right)}{d\tau_t}, &
    \frac{d\upsilon_{1,t}^0}{d\tau_t} &= \beta_{t,0}\left(\tilde{c}_{2,t+1}^0\right)^{1-1/\rho}\frac{d\ln\left(\tilde{c}_{2,t+1}^0\right)}{d\tau_t}, \\
    \frac{d\upsilon_{2,t}^i}{d\tau_t} &= \left(c_{2,t}^i\right)^{1-1/\rho}\frac{d\ln\left(c_{2,t}^i\right)}{d\tau_t}, &
    \frac{d\upsilon_{2,t}^0}{d\tau_t} &= \left(\tilde{c}_{2,t}^0\right)^{1-1/\rho}\frac{d\ln\left(\tilde{c}_{2,t}^0\right)}{d\tau_t}.
\end{align}
\end{subequations}
The young informal term uses that $\tilde{c}_{1,t}^0$ depends on $s_{t-1}$ only (see \eqref{\refmodeleq:EE:c0}), so its own derivative vanishes and only the $\tilde{c}_{2,t+1}^0$ channel survives.

For the young formal households, the discount factor $B_{t+1}^i$ is itself a function of $\tau_t$ through $s_t$, so the factor $(1+B_{t+1}^i)$ in $\upsilon_{1,t}^i$ cannot be held fixed while differentiating. It is convenient to absorb it into the consumption level, defining
\begin{align}\label{\refeq:hatc1i}
    \hat{c}_{1,t}^i \equiv \left(1+B_{t+1}^i\right)^{\frac{1}{1-1/\rho}}\tilde{c}_{1,t}^i \qquad\Longrightarrow\qquad \upsilon_{1,t}^i = \frac{\left(\hat{c}_{1,t}^i\right)^{1-1/\rho}}{1-1/\rho},
\end{align}
so that a single numerical derivative of $\ln(\hat{c}_{1,t}^i)$ picks up both channels at once.\footnote{In implementation $\hat{c}_{1,t}^i$ is never formed as a level. Only two things are ever needed from it, and both avoid the $1/(1-1/\rho)$ power: $\left(\hat{c}_{1,t}^i\right)^{1-1/\rho} = \left(1+B_{t+1}^i\right)\left(\tilde{c}_{1,t}^i\right)^{1-1/\rho}$ and $\ln\left(\hat{c}_{1,t}^i\right) = \ln\left(1+B_{t+1}^i\right)/(1-1/\rho)+\ln\left(\tilde{c}_{1,t}^i\right)$. This matters as $\rho\rightarrow 1$: the literal definition raises $1+B_{t+1}^i$ to a power that diverges, overflowing in double precision well before $\rho$ reaches $1$ (at $\rho = 1.001$ the exponent is already $1001$), whereas neither expression above does.}

Three of the log-derivatives in \eqref{\refeq:PEEcrraTerms} --- those of $h_t$, $\hat{c}_{1,t}^i$ and $\tilde{c}_{2,t+1}^0$ --- are approximated numerically along the $\tau_t$ dimension of the solution grid, since each depends on $\tau_t$ through the continuation policies and has no closed form. The remaining two do \emph{not} follow that route:
\begin{itemize}
    \item $d\ln(c_{2,t}^i)/d\tau_t$ \textbf{must} be taken from its closed form, given in \eqref{\refmodeleq:PEE}. Differentiating it numerically along the grid would be wrong, not merely imprecise: $s_{t-1,i}/s_{t-1}$ varies along that grid, and the policy maker takes it as predetermined, so the grid derivative would include a channel that does not belong in the first order condition.
    \item $d\ln(\tilde{c}_{2,t}^0)/d\tau_t$ \emph{may} be taken either way; we use its closed form for consistency with the above.
\end{itemize}
Both closed forms are exactly the expressions already used in the log case, \eqref{\refmodeleq:PEELOG}, with one substitution: the analytical $\partial\ln(\Theta_{h,t})/\partial\tau_t$ is replaced by the numerically obtained $d\ln(h_t)/d\tau_t$. Since $h_t = \Theta_{h,t}(s_{t-1}/\nu_t)^{\alpha\xi/(1+\alpha\xi)}$ and $s_{t-1}$ is held fixed under $\partial/\partial\tau_t$, the two coincide, and the log and CRRA cases share one formula rather than two.

\smalltitle{The terminal period, and the log limit.}
Under CRRA the terminal period needs no numerical differentiation at all: by the chain rule, each term of \eqref{\refmodeleq:terminalPEECRRA} is a consumption level raised to $1-1/\rho$ times the \emph{same} log-derivative already available in closed form from the log case. What is genuinely new at $t=T$ is that $B_T^i$ is no longer the primitive $\beta_i$, so the terminal problem is $z_T(\tau_T, s_{T-1})$ --- state-dependent, unlike its log counterpart. Two consequences are worth recording. First, the recursion at $t<T$ divides by $1-1/\rho$ (in \eqref{\refeq:PEEcrraTerms} and \eqref{\refeq:hatc1i}) and therefore refuses $\rho = 1$ outright, directing the caller to \cref{\refsec:PEELOG} rather than approximating the limit. Second, the terminal period is well defined at $\rho = 1$, where $\bm{B}_T$ collapses to $\bm{\beta}$ and every level factor to unity; it then coincides \emph{exactly} with the terminal period of the log solver, state-independence included. Since the two solvers share almost no code path, this identity --- and more generally the convergence of the CRRA solution to the log solution as $\rho\rightarrow 1$ --- is the sharpest available cross-check on both implementations.

\begin{alg}[Backward grid search over the savings level]
\label{\refalg:CRRA:grid}
Fix $\bm{\theta},\bm{\epsilon}$ and the grids $\mathcal{T}$, $\mathcal{S}$, $\mathcal{S}'$. The recursion identifies a sequence of policy functions $\tau_t(s_{t-1})$ together with the state policy functions $s_t(\tau_t, s_{t-1})$, iterating backwards from the terminal period $t=T$.

\smallskip\noindent\emph{Terminal period.} Every object entering $z_T$ is closed form in $(\tau_T, s_{T-1})$. On the full grid $\mathcal{S}\tilde{\mathcal{T}}$, compute and store: $d \ln(h_t)/d \tau_t$ (closed form for $t=T$), $\Theta_{h,t}$ (closed form for $t=T$), $h_t$, $B_t^i$, $\Gamma_{s,t-1}$, $s_{t-1,i}/s_{t-1}$, $c_{2,t}^i$, $\tilde{c}_{2,t}^0$, and $\tilde{c}_{1,t}^i$ (closed form for $t=T$). Assemble the first order condition $z_T$ from these, and apply the selection rule of \cref{\refsec:RobustRoot} along $\tau$ for each state --- boundary checks, direction filter and objective integration included --- to obtain $\tau_T(s_{T-1})$ on $\mathcal{S}$. Finally, re-evaluate the auxiliary objects at the selected taxes and construct the interpolants $\tau^T(\cdot)$ and $h^T(\cdot)$ that the previous period will call.

\smallskip\noindent\emph{Recursion, $t<T$.} Given the interpolants $\tau^{t+1}(\cdot)$ and $h^{t+1}(\cdot)$, each period proceeds in four steps.

\begin{enumerate}
    \item \emph{State approximation.} On the Cartesian grid $\tilde{\mathcal{T}}\times\mathcal{S}'$ evaluate the forward pass \eqref{\refeq:stateApprox}, retaining $\Theta_{s,t}$ (and $\Theta_{h,t}$, $\Gamma_{s,t}$, $\tau^{t+1}$, $h^{t+1}$, $\bm{B}_{t+1}$ for later use). Form the residual \eqref{\refeq:stateResidual} for every $s_{t-1}\in\mathcal{S}$ by broadcasting, and locate its root along the $s_t$ dimension, giving $s_t(\tau_t, s_{t-1})$ on $\tilde{\mathcal{T}}\times\mathcal{S}$ together with the feasibility mask. This is the only step that touches $\mathcal{S}'$, and the only one that calls the continuation interpolants.

    \item \emph{Current-period objects.} With $s_t$ resolved, evaluate on $\tilde{\mathcal{T}}\times\mathcal{S}$: $h_t\left(\Theta_{h,t}, s_{t-1}\right)$, then $R_t$, $\bm{B}_t$, $\Gamma_{s,t-1}\left(\bm{B}_t, \tau_t\right)$ and $s_{t-1,i}/s_{t-1}$ --- the last two evaluated at the previous period's parameter vintage, exactly as in the log case. Then the consumption levels entering \eqref{\refeq:PEEcrraTerms}: $\hat{c}_{1,t}^i$ from \eqref{\refeq:hatc1i}, $c_{2,t}^i$, $\tilde{c}_{2,t}^0$, and $\tilde{c}_{2,t+1}^0$ (a function of $s_t$, $\tau^{t+1}(s_t)$ and $h^{t+1}(s_t)$).

    \item \emph{Derivatives and the first order condition.} Approximate $d\ln(h_t)/d\tau_t$, $d\ln(\hat{c}_{1,t}^i)/d\tau_t$ and $d\ln(\tilde{c}_{2,t+1}^0)/d\tau_t$ numerically along $\tau_t$, separately for each $s_{t-1}$; take $d\ln(c_{2,t}^i)/d\tau_t$ and $d\ln(\tilde{c}_{2,t}^0)/d\tau_t$ from their closed forms as described above. Assemble $z_t$ on $\tilde{\mathcal{T}}\times\mathcal{S}$ via \eqref{\refeq:PEEcrraTerms} and the political weights of \eqref{\refeq:LOG:fast}, leaving infeasible cells undefined.

    \item \emph{Selection and reporting.} Apply the selection rule \eqref{\refeq:candidates} along $\tau_t$ for each $s_{t-1}$, skipping infeasible cells, to obtain $\tau_t(s_{t-1})$. Re-evaluate steps 1--2 at the selected taxes to obtain the solution dictionary on $\mathcal{S}$, and construct the interpolants $\tau_t(s_{t-1})$ and $h_t(s_{t-1})$ that the previous period will call. A light smoothing of the selected $\tau_t(s_{t-1})$ before interpolation removes small kinks inherited from the interpolated continuation policies; the smoother's interior knots are fixed in advance (at every $m$-th valid node) rather than selected from the data, so that the smoothed profile is a \emph{linear} map of the profile it smooths --- this keeps the composed map from parameters to solution continuous, which matters because the calibration of \cref{\refsec:calibration} differentiates through it.
\end{enumerate}

Steps 1--4 are stated for the whole grid $\mathcal{S}$ at once rather than for one $s_{t-1}$ at a time. Nothing forces this --- the states are independent and could equally be looped over --- but the grid evaluations dominate the cost only through the number of \emph{calls}, not the number of gridpoints, so handling all states in one pass is materially cheaper than looping, at the price of carrying the feasibility mask explicitly rather than subsetting the grid per state.
\end{alg}
